\documentclass[17pt]{flash}
\usepackage{texmaths-tikz}
\startdatakeys
    \BO{2026}
    \cycle{4}
    \competencesDuSocle{RA}
    \automatismes{C4A1cd}
\enddatakeys
\begin{document}

\flashvisibletimertrue
\FlashSetItemSep{40pt}
\FlashSetVtop{-1em}
\FlashSetTitle{Lire une fraction sur une droite graduée}
\FlashShowConsignes{%
    \item Dans chaque cas, écrire l'abscisse du point \mathspoint{M}. 
}
\flashlist{
    \flashitem{45}{
        \begin{tikzpicture}[scale=1]
            \coordinate(O)at(0,0);
            \node at([shift={(-90:0.4)}]O){0};
            \coordinate(A)at(4,0);
            \node at([shift={(-90:0.4)}]A){1};
            \foreach\x in{1,2,3}{
                \draw(\x,-0.1)--(\x,0.1);
            }
            \draw[line width=1pt](0,-0.15)--(0,0.15);
            \draw[line width=1pt](4,-0.15)--(4,0.15);
            \tkzDrawLine[add=0.1 and 0.25, -Stealth](O,A)
            \coordinate(B)at(1,0);
            \node at([shift={(90:0.4)}]B){\mathspoint{M}};
            \draw[line width=1.2pt, color=\markscolor](1,-0.15)--(1,0.15);
        \end{tikzpicture}
    }
    \flashitem{45}{
        \begin{tikzpicture}[scale=0.4]
            \coordinate(O)at(0,0);
            \node at([shift={(-90:0.8)}]O){\small 0};
            \coordinate(A)at(10,0);
            \node at([shift={(-90:0.8)}]A){\small 1};
            \foreach\x in{1,...,10}{
                \draw[line width=0.8pt](\x,-0.1)--(\x,0.1);
            }
            \draw[line width=1pt](0,-0.15)--(0,0.15);
            \draw[line width=1pt](10,-0.15)--(10,0.15);
            \tkzDrawLine[add=0.1 and 0.25, -Stealth](O,A)
            \coordinate(B)at(7,0);
            \node at([shift={(90:0.8)}]B){\small\mathspoint{M}};
            \draw[line width=1.2pt, color=\markscolor](7,-0.15)--(7,0.15);
        \end{tikzpicture}
    }
    \flashitem{45}{
        \begin{tikzpicture}[scale=0.5]
            \coordinate(O)at(0,0);
            \node at([shift={(-90:0.8)}]O){\small 0};
            \coordinate(A)at(6,0);
            \node at([shift={(-90:0.8)}]A){\small 2};
            \foreach\x in{1,...,5}{
                \draw[line width=0.8pt](\x,-0.1)--(\x,0.1);
            }
            \draw[line width=1pt](0,-0.15)--(0,0.15);
            \draw[line width=1pt](6,-0.15)--(6,0.15);
            \tkzDrawLine[add=0.1 and 0.25, -Stealth](O,A)
            \coordinate(B)at(4,0);
            \node at([shift={(90:0.8)}]B){\small\mathspoint{M}};
            \draw[line width=1.2pt, color=\markscolor](4,-0.15)--(4,0.15);
        \end{tikzpicture}
    }
}

\end{document}
